% Partie : théorie de la démonstration
% Chapitre : calcul des séquents
% Section : interprétation des règles

\documentclass[../../../include/open-logic-section]{subfiles}

\begin{document}

\olfileid[fr]{pt}{seq}{irl}

\olsection{Interprétation des règles}

Rappelons l'interprétation d'un séquent :
$\Gamma \Sequent \Delta$ est vrai dans une
structure~$\Struct{M}$ si au moins une formule
$!A \in \Gamma$ est fausse ou si au moins une formule
$!A \in \Delta$ est vraie.
Les règles propositionnelles de~\Log{G3c}
expriment les conditions de vérité de leurs
formules principales. Considérons \RightR{\lif}.
Sa principale $!A \lif !B$ est vraie si $!A$
est fausse ou si~$!B$ est vraie.
Le séquent $\ \Sequent !A \lif !B$ est donc
vrai exactement quand $!A \Sequent !B$ est vrai.
En ajoutant les contextes $\Gamma$ et $\Delta$
aux antécédents et aux conséquents, on obtient
la conclusion et la prémisse de \RightR{\lif}.

Pour~\LeftR{\lif}, la formule $!A \lif !B$
est fausse exactement quand $!A$ est vraie
et $!B$ est fausse.
Ainsi, $!A \lif !B \Sequent \ $ est vrai
exactement quand $\ \Sequent !A$ et
$!B \Sequent \ $ sont tous deux vrais.
La conclusion de \LeftR{\lif} est équivalente
à la conjonction de ses prémisses.
Cette équivalence rend ces règles correctes :
si toutes les prémisses sont vraies, la conclusion
est vraie. Elle intervient aussi dans la preuve
de complétude du calcul.

Les conditions de vérité de tout connecteur
défini par une fonction de vérité se décrivent de cette façon par
des séquents. Toute fonction de vérité classique
se représente par une formule en forme normale
conjonctive : une conjonction de disjonctions
de formules atomiques ou de leurs négations.
Prenons, par exemple, $!A \oplus !B$, vraie
exactement quand une seule des formules $!A$
et $!B$ est vraie : le \emph{ou exclusif}.
Elle est vraie exactement quand
$(!A \lor !B) \land (\lnot !A \lor \lnot !B)$
est vraie, et fausse exactement quand
$(\lnot !A \lor !B) \land (!A \lor \lnot !B)$
est vraie. On peut donc introduire les règles pour $\oplus$ :
\begin{align*}
&
\Axiom$\Gamma \fCenter \Delta, !A, !B$
\Axiom$!A, !B, \Gamma \fCenter \Delta$
\RightLabel{\RightR{\oplus}}
\BinaryInf$\Gamma \fCenter \Delta, !A \oplus !B$
\DisplayProof
\\
& \Axiom$!A, \Gamma \fCenter \Delta, !B$
\Axiom$!B, \Gamma \fCenter \Delta, !A$
\RightLabel{\LeftR{\oplus}}
\BinaryInf$!A \oplus !B, \Gamma \fCenter \Delta$
\DisplayProof
\end{align*}
qui donnent une présentation correcte et complète
de ce connecteur : chaque conclusion est vraie
exactement quand toutes ses prémisses sont vraies.

\begin{prob}
Supposons $!A \downarrow !B$ vraie exactement
quand ni $!A$ ni~$!B$ n'est vraie : c'est
le connecteur \emph{non-ou}, ou NOR.
Quelles seraient ses règles dans \Log{G3c} ?
Trouver deux séquents simultanément vrais
exactement quand $!A \downarrow !B$ est vraie :
ils donnent la règle \RightR{\downarrow}.
Puis trouver un séquent vrai exactement quand
$!A \downarrow !B$ est fausse :
il donne la règle \LeftR{\downarrow}.
\end{prob}

Dans \Log{G3c}, chaque connecteur ou quantificateur
a une règle gauche et une règle droite.
Dans \Log{G1c}, on a cependant deux règles
\LeftR{\land} et deux règles \RightR{\lor}.
Le système original de Gentzen, \Log{LK},
sur lequel \Log{G1c} est modelé, avait une
variante~\Log{LJ} pour la logique intuitionniste.
Cette variante a des conséquents contenant
au plus une formule et des règles adaptées
à cette restriction. La règle \RightR{\lor}
de~\Log{G3c} ne convient pas, puisque sa prémisse
contient $!A$ et~$!B$ dans le conséquent.
Gentzen a donc employé deux règles
\RightR{\lor} dans \Log{LJ}, et les mêmes
dans~\Log{LK}.
De même, \Log{G1i} a au plus une formule
dans chaque conséquent et adapte les règles
de~\Log{G1c}, notamment l'implication gauche.
La variante intuitionniste de \Log{G3c}
modifie aussi certaines règles ; elle emploie
elle aussi deux règles \RightR{\lor}.

La correction des règles structurelles de~\Log{G1c}
est facile à vérifier. Si
$\Gamma \Sequent \Delta$ contient une formule
fausse à gauche ou une formule vraie à droite,
$\Gamma \Sequent \Delta, !A$ conserve cette
même formule, quelle que soit la valeur de $!A$.
Si la prémisse de~\RightR{\Weakening} est
vraie, sa conclusion l'est donc aussi.
La réciproque ne vaut pas : la conclusion
peut être vraie seulement parce que $!A$
est vraie, sans que la prémisse soit vraie.

À quoi servent les règles structurelles ?
La contraction, \RightR{\Contraction}
ou \LeftR{\Contraction}, permet notamment
de construire une preuve de
$\ \Sequent !A \lor \lnot !A$.
Partant de $!A \Sequent !A$, on obtient
$\Sequent !A, \lnot !A$ par~\RightR{\lnot}.
On applique ensuite successivement les deux
versions de \RightR{\lor}, en remplaçant
d'abord $\lnot !A$, puis~$!A$, par
$!A \lor \lnot !A$.
Il reste le séquent
$\ \Sequent !A \lor \lnot !A, !A \lor \lnot !A$.
Pour finir cet arbre dans~\Log{G1c} par
$\Sequent !A \lor \lnot !A$, on contracte
les deux occurrences :
\[
\Axiom$!A \fCenter !A$
\RightLabel{\RightR{\lnot}}
\UnaryInf$\fCenter !A, \lnot !A$
\RightLabel{\RightR{\lor}}
\UnaryInf$\fCenter !A, !A \lor \lnot !A$
\RightLabel{\RightR{\lor}}
\UnaryInf$\fCenter !A \lor \lnot !A, !A \lor \lnot !A$
\RightLabel{\RightR{\Contraction}}
\UnaryInf$\fCenter !A \lor \lnot !A$
\DisplayProof
\]

Certains séquents valides ne se prouvent pas
dans~\Log{G1c} sans affaiblissement ni contraction.
Prenons $!A$ et $!B$ comme deux variables
propositionnelles distinctes, et considérons
$!A \Sequent !B \lif !A$.
En l'absence de règles structurelles, la dernière
inférence devrait être \RightR{\lif} :
seule $!B \lif !A$ pourrait être principale.
Sa prémisse serait $!B, !A \Sequent !A$,
qui n'est pas un axiome et ne peut être
la conclusion d'aucune règle logique,
puisque toutes ses formules sont atomiques.
Les règles structurelles sont donc indispensables
dans cet exemple. Avec \LeftR{\Weakening},
on obtient la prémisse à partir de
$!A \Sequent !A$ :
\[
\Axiom$!A \fCenter !A$
\RightLabel{\LeftR{\Weakening}}
\UnaryInf$!B, !A \fCenter !A$
\RightLabel{\RightR{\lif}}
\UnaryInf$!A \fCenter !B \lif !A$
\DisplayProof
\]

Dans \Log{G3c}, affaiblissement et contraction
ne sont pas des règles primitives : ils sont
admissibles. L'affaiblissement est déjà incorporé
dans les contextes des axiomes.
Les règles à deux prémisses de \Log{G1c}
et de \Log{G3c} partagent leurs contextes
$\Gamma$ et $\Delta$, ce qui évite de réunir
deux copies de chaque contexte dans la conclusion.
Pour les règles à une prémisse, la contraction
intervient notamment avec \RightR{\lor},
\LeftR{\land}, \RightR{\lexists}
et~\LeftR{\lforall} dans la présentation G1c.
Les versions G3c évitent ces contractions explicites.
Le calcul \Log{G2c}, présenté dans le \olref{tab:G2c},
a au contraire des contextes indépendants
dans ses règles à deux prémisses.
La contraction y joue donc un rôle encore
plus important que dans~\Log{G1c}.

\subfile{rules-G2c}

Pour les règles à constante propre, la preuve de
correction fait varier l'interprétation de cette constante.
Par exemple, si la conclusion de \RightR{\lforall}
est fausse dans une structure, un élément du domaine
réfute la formule universelle. On donne cet élément
pour valeur à la constante propre, absente de la
conclusion. La prémisse est alors fausse dans la
%
%
structure ainsi obtenue.

\end{document}
